\section{结论}
本文研究了随机触发远程估计中可信预测反馈与可篡改载荷并存时，指定记录能够为真实估计损伤提供何种约束。精确反馈使隐藏估计偏差进入下一可信事件的条件概率；目标协方差标签校准不同事件历史的健康基准。基于这两点，本文对全部允许的有界因果攻击建立了固定有限时域的双预算损伤界，并由合法构造得到匹配阶数。

标量量化扩展刻画了该保证的适用边界：当量化后的预测中心参与实际触发、且审计记录限于指定遥测时，内部预测可在同一量化格内变化而不改变该记录，因而零TV仍可伴随正损伤。该结果针对量化协议和受限记录的组合，不表示攻击对完整接收日志不可检测。增强审计能区分本文给出的见证，但本文没有证明它排除全部零TV正损伤策略。

数值结果用于解释损伤分解、解析双预算界和量化见证。精确反馈样本没有观察到事件分叉，不能据此判断总体TV为零；量化样本只覆盖精度主导区域，不能替代全域定理。结论限于全状态含噪观测、可逆状态矩阵、固定有限目标时刻及明示的信息结构。证书常数可能保守；无饱和量化不代表固定有限比特实现；总体TV界也不是有限样本在线认证程序。
